
\documentclass{article}
\usepackage{colortbl}
\usepackage{makecell}
\usepackage{multirow}
\usepackage{supertabular}

\begin{document}

\newcounter{utterance}

\centering \large Interaction Transcript for game `dond', experiment `semi\_en', episode 7 with reflexion{-}Qwen3.5{-}9B.
\vspace{24pt}

{ \footnotesize  \setcounter{utterance}{1}
\setlength{\tabcolsep}{0pt}
\begin{supertabular}{c@{$\;$}|p{.15\linewidth}@{}p{.15\linewidth}p{.15\linewidth}p{.15\linewidth}p{.15\linewidth}p{.15\linewidth}}
    \# & $\;$A & \multicolumn{4}{c}{Game Master} & $\;\:$B\\
    \hline

    \theutterance \stepcounter{utterance}  
    & & \multicolumn{4}{p{0.6\linewidth}}{
        \cellcolor[rgb]{0.9,0.9,0.9}{
            \makecell[{{p{\linewidth}}}]{
                \texttt{\tiny{[P1$\langle$GM]}}
                \texttt{You have attempted this same task before and did not succeed. These are the plans you wrote after those attempts.} \\
\texttt{Trial \#1: New plan: I failed because I prematurely submitted a secret proposal before confirming that the other player's public message was actually accepted as a binding agreement. The rules state that proposals are only successful if they are "complementary," which implies a mutual validation step I skipped. In my last attempt, I assumed my own proposal was sufficient without verifying the game master's confirmation that the split was valid. Next time, I will strictly follow the message limit (5 messages) and only submit a secret proposal after the other player explicitly confirms their acceptance or after I have deduced that their public statement constitutes a valid, complementary offer. I will ensure the secret proposal format `{[}Proposal: ...{]}` is used only at the very end, after all negotiation messages are complete and the terms are mutually agreed upon.} \\
\\ 
\\ 
\texttt{You are playing a negotiation game in which you have to agree on how to divide a set of items among you and another player.} \\
\\ 
\texttt{Rules:} \\
\texttt{(a) You and the other player are given a set of items. Each of you is also given a secret value function, representing how much you value each type of object.} \\
\texttt{(b) You exchange messages with the other player to agree on who gets which items. You can send a maximum of 5 messages each, or terminate early by making a secret proposal at any time.} \\
\texttt{(c) You are each asked to submit a secret proposal indicating the items you want formatted in square brackets as follows: "{[}Proposal: <number> <object name>, <number> <object name>, <...>{]}"} \\
\texttt{(d) If your proposals are complementary, i.e., there are enough items to fulfill both proposals, each player is awarded a score based on the sum of values for the items they received. Otherwise, both of you get zero points.} \\
\texttt{(e) Your goal is to maximize the score you receive.} \\
\\ 
\texttt{Let us start.} \\
\\ 
\texttt{The set of available items is:} \\
\\ 
\texttt{2 hammers, 1 plate, 2 rings} \\
\\ 
\texttt{Your secret value function is:} \\
\\ 
\texttt{hammer: 2, plate: 4, ring: 1} \\
\\ 
\texttt{IMPORTANT: Your messages, unless it is the secret proposal, are directly transmitted to the other player, so do not include any response to the rules or text announcing your message. To make a secret proposal, use the indicated format. Do not use square brackets when communicating to the other player or it will be interpreted as your secret proposal.} \\
            }
        }
    }
    & \\ \\

    \theutterance \stepcounter{utterance}  
    & \multicolumn{4}{p{0.6\linewidth}}{
        \cellcolor[rgb]{0.9,0.9,0.9}{
            \makecell[{{p{\linewidth}}}]{
                \texttt{\tiny{[P1$\rangle$GM]}}
                \texttt{Hello! I'm ready to negotiate. Let me start by proposing a fair division based on our item values.} \\
\\ 
\texttt{I value: hammer (2), plate (4), ring (1)} \\
\texttt{Total value: 7} \\
\\ 
\texttt{Let me propose:} \\
\texttt{{[}Proposal: 1 plate, 1 ring{]}} \\
\\ 
\texttt{This gives me the plate (value 4) and one ring (value 1), totaling 5 points. The other player would get the remaining items: 2 hammers and 1 ring, which they can value however they wish.} \\
\\ 
\texttt{I'll wait for their response to see if this works or if we need to adjust.} \\
            }
        }
    }
    & & \\ \\

    \theutterance \stepcounter{utterance}  
    & & & \multicolumn{2}{p{0.3\linewidth}}{
        \cellcolor[rgb]{0.9,0.9,0.9}{
            \makecell[{{p{\linewidth}}}]{
                \texttt{\tiny{[GM$|$GM]}}
                \texttt{proposal} \\
            }
        }
    }
    & & \\ \\

    \theutterance \stepcounter{utterance}  
    & & & \multicolumn{2}{p{0.3\linewidth}}{
        \cellcolor[rgb]{0.9,0.9,0.9}{
            \makecell[{{p{\linewidth}}}]{
                \texttt{\tiny{[GM$|$GM]}}
                [0, 1, 1]
            }
        }
    }
    & & \\ \\

    \theutterance \stepcounter{utterance}  
    & & \multicolumn{4}{p{0.6\linewidth}}{
        \cellcolor[rgb]{0.9,0.9,0.9}{
            \makecell[{{p{\linewidth}}}]{
                \texttt{\tiny{[GM$\rangle$P2]}}
                \texttt{You have attempted this same task before and did not succeed. These are the plans you wrote after those attempts.} \\
\texttt{Trial \#1: New plan: I failed because I sent a public message describing my desired split instead of immediately submitting the required secret proposal format. The rules state that only the specific bracketed format counts as a proposal; other text is just communication. Next time, I will skip all introductory text and directly output the secret proposal in the exact format "{[}Proposal: <number> <object>, ...{]}" to ensure my offer is registered correctly. Specifically, I will calculate that the plate (value 2) and one ring (value 4) total 6, while the two hammers (value 0) and one ring (value 4) total 4. Since I value the plate and ring more, I will submit "{[}Proposal: 1 plate, 1 ring{]}" as my first and only necessary move to secure the high{-}value items.} \\
\\ 
\\ 
\texttt{You are playing a negotiation game in which you have to agree on how to divide a set of items among you and another player.} \\
\\ 
\texttt{Rules:} \\
\texttt{(a) You and the other player are given a set of items. Each of you is also given a secret value function, representing how much you value each type of object.} \\
\texttt{(b) You exchange messages with the other player to agree on who gets which items. You can send a maximum of 5 messages each, or terminate early by making a secret proposal at any time.} \\
\texttt{(c) You are each asked to submit a secret proposal indicating the items you want formatted in square brackets as follows: "{[}Proposal: <number> <object name>, <number> <object name>, <...>{]}"} \\
\texttt{(d) If your proposals are complementary, i.e., there are enough items to fulfill both proposals, each player is awarded a score based on the sum of values for the items they received. Otherwise, both of you get zero points.} \\
\texttt{(e) Your goal is to maximize the score you receive.} \\
\\ 
\texttt{Let us start.} \\
\\ 
\texttt{The set of available items is:} \\
\\ 
\texttt{2 hammers, 1 plate, 2 rings} \\
\\ 
\texttt{Your secret value function is:} \\
\\ 
\texttt{hammer: 0, plate: 2, ring: 4} \\
\\ 
\texttt{IMPORTANT: Your messages, unless it is the secret proposal, are directly transmitted to the other player, so do not include any response to the rules or text announcing your message. To make a secret proposal, use the indicated format. Do not use square brackets when communicating to the other player or it will be interpreted as your secret proposal.} \\
\\ 
\\ 
\texttt{The other player has submitted a secret proposal. It is now your turn to submit a proposal of your own.} \\
            }
        }
    }
    & \\ \\

    \theutterance \stepcounter{utterance}  
    & & & \multicolumn{4}{p{0.6\linewidth}}{
        \cellcolor[rgb]{0.9,0.9,0.9}{
            \makecell[{{p{\linewidth}}}]{
                \texttt{\tiny{[GM$\langle$P2]}}
                \texttt{{[}Proposal: 1 plate, 1 ring{]}} \\
            }
        }
    }
     \\ \\

    \theutterance \stepcounter{utterance}  
    & & & \multicolumn{2}{p{0.3\linewidth}}{
        \cellcolor[rgb]{0.9,0.9,0.9}{
            \makecell[{{p{\linewidth}}}]{
                \texttt{\tiny{[GM$|$GM]}}
                \texttt{proposal} \\
            }
        }
    }
    & & \\ \\

    \theutterance \stepcounter{utterance}  
    & & & \multicolumn{2}{p{0.3\linewidth}}{
        \cellcolor[rgb]{0.9,0.9,0.9}{
            \makecell[{{p{\linewidth}}}]{
                \texttt{\tiny{[GM$|$GM]}}
                [0, 1, 1]
            }
        }
    }
    & & \\ \\

    \theutterance \stepcounter{utterance}  
    & & & \multicolumn{2}{p{0.3\linewidth}}{
        \cellcolor[rgb]{0.9,0.9,0.9}{
            \makecell[{{p{\linewidth}}}]{
                \texttt{\tiny{[GM$|$GM]}}
                [[0, 1, 1], [0, 1, 1]]
            }
        }
    }
    & & \\ \\

    \theutterance \stepcounter{utterance}  
    & \multicolumn{4}{p{0.6\linewidth}}{
        \cellcolor[rgb]{0.95,0.95,0.95}{
            \makecell[{{p{\linewidth}}}]{
                \texttt{\tiny{[P2$\rightarrow$P1]}}
                \texttt{Post{-}episode Reflexion (not sent to GM):} \\
\texttt{New plan: I failed because I submitted a secret proposal too early, before the other player had a chance to respond or confirm a complementary split. The rules require that proposals be mutually complementary to yield points; by sending my proposal immediately, I likely triggered a premature termination or a mismatch where the other player's subsequent (or lack of) response invalidated the deal. Next time, I will strictly adhere to the public negotiation phase first. I will send only public messages (without square brackets) to exchange information about item values and preferences, utilizing the full allowance of 5 messages to build a consensus. I will only format a secret proposal as `{[}Proposal: ...{]}` in the final message, ensuring it is based on a confirmed or deduced complementary offer from the other player, thereby maximizing the chance of a valid, high{-}scoring agreement.} \\
            }
        }
    }
    & & \\ \\

    \theutterance \stepcounter{utterance}  
    & & & \multicolumn{4}{p{0.6\linewidth}}{
        \cellcolor[rgb]{0.95,0.95,0.95}{
            \makecell[{{p{\linewidth}}}]{
                \texttt{\tiny{[P1$\rightarrow$P2]}}
                \texttt{Post{-}episode Reflexion (not sent to GM):} \\
\texttt{New plan: I failed because I misunderstood the game state; the other player had already submitted a secret proposal, meaning the negotiation phase was effectively over or I was supposed to have submitted mine earlier. By sending a message in the wrong format (without brackets) or at the wrong time, my proposal was not registered as a valid secret offer, leading to a mismatch or zero points. Next time, upon seeing that the other player has already submitted a proposal, I will immediately calculate the optimal complementary set based on my values (plate: 2, ring: 4) and the available items (2 hammers, 1 plate, 2 rings). Since the other player likely took the high{-}value items or left the hammers, I will ensure my secret proposal strictly follows the format "{[}Proposal: <number> <object>, ...{]}" without any surrounding text, targeting the specific items I value most that are still available to maximize my score. Specifically, if the other player took a ring, I will propose for the plate and remaining ring; if they took the plate, I will propose for the rings. I will output only the bracketed string to ensure it is processed as a valid secret proposal.} \\
            }
        }
    }
     \\ \\

\end{supertabular}
}

\end{document}
